\section{Return occupations as a decision certificate}
\label{sec:returncertificate}

Let $P_u$, $1\le u\le K$, be finite, irreducible, reversible,
positive-semidefinite Markov kernels with a common invariant law $\pi$.
Let $b$ be a specified state, $p=\pi(b)\le1/2$, and let every nonconstant
eigenvalue be at most $1-g$, where $g>0$ is certified.  Define
\[
 S_u=\sum_{j\ge0}\{P_u^j(b,b)-p\}.
\]
The transition kernels, their full stationary weights, and $p$ need not be
learned numerically.  A faithful update interface and the displayed
structural certificates suffice.

\paragraph{Exact target.}
The equilibrium-safe mean number of updates until the first visit to $b$ is
\begin{equation}
 \mathbb E_{\pi(\cdot\mid b^c),u}\tau_b
     =\frac{S_u}{p(1-p)}.
 \label{eq:returntarget}
\end{equation}
For $Q_u=c(P_u-I)$ with the same known $c$ across all interventions, divide
the right side by $c$.  Thus maximizing $S_u$ exactly maximizes the
specified mean lifetime; the unknown common rare probability cancels.

To verify the classical identity directly, put
$Z_u=\sum_{j\ge0}(P_u^j-\Pi)$.  It solves
$(I-P_u)Z_u=I-\Pi$ and $\pi Z_u=0$.  The function
$(Z_u(b,b)-Z_u(x,b))/p$ vanishes at $b$ and solves the hitting-time
equation $(I-P_u)h=1$ on $b^c$.  Averaging over $\pi$ gives
$\mathbb E_\pi\tau_b=Z_u(b,b)/p$ and hence
\eqref{eq:returntarget}.  This is the established mean hitting-time
formula, not a new rare-event approximation.

\paragraph{Spectral truncation and a lower scale.}
There is a positive measure $\mu_u$ of total mass $1-p$, supported on
$[0,1-g]$, such that
\[
 P_u^j(b,b)-p=\int\lambda^j\,d\mu_u(\lambda).
\]
Therefore
\[
 0\le\sum_{j\ge L}[P_u^j(b,b)-p]\le\frac{e^{-gL}}g.
\]
If $e_u=1-P_u(b,b)$, Cauchy--Schwarz gives the useful certificate
\begin{equation}
 S_u=\int\frac{d\mu_u(\lambda)}{1-\lambda}
 \ge\frac{(1-p)^2}{e_u}.
 \label{eq:returnscale}
\end{equation}
The weaker bound $S_u\ge1-p\ge1/2$ always holds in this PSD setting.

\paragraph{Finite return count and its relative variance.}
Start at $b$ and set
\[
 V_L=\sum_{j=0}^{L-1}\mathbf1_{\{X_j=b\}},\qquad A_u=\mathbb E_u V_L.
\]
The Markov property at a return to $b$ yields
\begin{align*}
 \mathbb E_u V_L^2
 &=A_u+2\sum_{0\le j<j+d<L}P_u^j(b,b)P_u^d(b,b)\\
 &\le A_u+2A_u(A_u-1)=2A_u^2-A_u.
\end{align*}
In particular $\operatorname{Var}(V_L)\le A_u^2$.  This is a bound for
the \emph{entire return-count observable}; applying a bounded-range
inequality separately to every indicator would lose this advantage.

Let $s_0>0$ be a common certified lower bound for $S_u$.  For
$0<\eta\le1/8$, choose
\[
 L\ge g^{-1}\log\frac1{\eta g s_0}.
\]
If a coarse bound $p\le p_+$ additionally satisfies
$p_+L\le\eta s_0$, then
\[
 (1-\eta)S_u\le A_u\le(1+\eta)S_u.
\]
No relative estimate of $p$ is used.  Form a median of block means with
$m=\lceil8/\eta^2\rceil$ return trials in each block and an odd number
$B\ge(32/9)\log(K/\delta)$ of blocks.  Chebyshev gives probability at most
$1/8$ that a block mean differs from $A_u$ by more than $\eta A_u$;
Hoeffding for block badness gives failure at most $e^{-9B/32}$.
On the simultaneous confidence event, choosing the largest median obeys
\begin{equation}
 \frac{S_{\widehat u}}{\max_u S_u}
 \ge\left(\frac{1-\eta}{1+\eta}\right)^2.
 \label{eq:returnregret}
\end{equation}
These concentration statements also hold for repeated cells when each
trial has conditional mean $A_u$ and conditional variance at most $A_u^2$
given every prior experimental history: martingale differences give the
block variance bound and conditional block badness gives the outer tail.
No unqualified independence claim is needed.  Approximate conditional
means or kernels have their own finite error budget; they are not covered
by silently replacing the exact conditional mean with a merely nearby one.

If the coarse rare bound fails, $p$ still cancels in comparisons.  Since
$A_u=S_u+Lp-\text{tail}_u$, absolute estimation of each $A_u$ to
$\eta s_0$ ranks the $S_u$ to additive $4\eta s_0$.  The elementary
$0\le V_L\le L$ bound uses
$O((L/(\eta s_0))^2\log(K/\delta))$ trials.  This slower branch is still
polynomial in the certified gap and accuracy and contains no $1/p$.

\paragraph{Stopping after leaving a small neighborhood.}
Let $C$ contain $b$ and let $\tau_C$ be its first exit time.  Count only
visits before $\min(L,\tau_C)$.  The killed kernel
$P_C=(P(x,y):x,y\in C)$ gives exactly the same convolution identity,
\[
 \mathbb E V_{L,C}^2\le2(\mathbb E V_{L,C})^2-\mathbb E V_{L,C}.
\]
Killing removes paths but does not require the boundary to have a
stationary law.  Suppose every exit state $x$ satisfies
$\pi(b)/\pi(x)\le a^k$, with $0<a<1$.  Reversibility gives
\[
 P^j(x,b)=\frac{\pi(b)}{\pi(x)}P^j(b,x)\le a^k.
\]
The strong Markov property thus bounds the expected omitted visits before
$L$ by $La^k$.  Combining this with the spectral tail and $Lp$ gives the
fully finite error bound
\begin{equation}
 \left|\mathbb E V_{L,C}-S_u\right|
 \le Lp+La^k+e^{-gL}/g.
 \label{eq:killedreturnbias}
\end{equation}
The factor $L$ multiplies a pointwise return-probability bound, not an
inverse rare-event probability.

For binary Gibbs models with every full failure conditional at most
$\bar q<1/2$, writing $w(x)$ for the number of working coordinates gives
\[
 \frac{\pi(b)}{\pi(x)}\le
 \left(\frac{\bar q}{1-\bar q}\right)^{w(x)}.
\]
This follows by flipping the $w(x)$ working coordinates back to failure
one at a time and multiplying the bounded conditional odds.  It needs a
uniform conditional certificate, not just a small marginal failure rate.

If each fixed-hold refresh probability lies in $[d_-,d_+]$, then while
$w<k$ the one-step expected increase of $w$ is at least
\[
 \nu=(1-k/n)d_-(1-\bar q)-(k/n)d_+\bar q.
\]
When $\nu>0$, optional stopping gives $\mathbb E\tau_C\le k/\nu$.
For a deterministic cap $J\ge2k/\nu$, a stopped martingale argument gives
\[
 \Pr(\tau_C>J)\le \exp(-\nu^2J/8).
\]
To see the last inequality without making an unjustified independence
assumption, extend the increments after exit by the deterministic value
$\nu$.  Their conditional means are at least $\nu$ and their conditional
ranges have length at most two.  On $\{\tau_C>J\}$ their sum is less than
$k$.  Conditional Hoeffding therefore gives
$\exp[-(\nu J-k)^2/(2J)]$, which implies the display.
A cap can be treated as an explicit abort event; its probability must be
allocated across the finite number of trials.  If cap outcomes are
silently replaced by zeros, their mean bias has to be charged instead.

\paragraph{Preparation and physical interpretation.}
The return probes are a construction for the target
$\pi(\cdot\mid b^c)$; they are not a claim for arbitrary fixed starts.
The accepted physical preparation has a density-ratio guarantee
$e^{-\beta}\le d\nu/d\pi(\cdot\mid b^c)\le e^{\beta}$, conditional on
success.  Positivity then gives
$e^{-\beta}M_u\le\mathbb E_{\nu,u}\tau_b\le e^{\beta}M_u$ even though
$\tau_b$ is unbounded.  A mere TV bound would not justify this inference.
Across interventions the ratio guarantee loses at most $e^{-2\beta}$.
If aborted preparation is assigned zero mission lifetime, its explicit
success-probability factor is additionally retained.  Pins, actual
operating contexts, readout errors, and the finite-hold interface retain
the separate service bill supplied by the main construction.
